% Section 5: Empirical Validation
% IEEE Computer — "Silent Decisions Are Type Errors"
% Authors: Soares et al., 2026
%
% OS-10 (COMSI-2026-04-0112, closes A2): every numeric cell below is
% traceable to a committed CSV under data/ (see scripts/summarize_sweep.py,
% which prints the exact median/CV values transcribed into the tables) or
% to a fresh `cargo bench` run of btv-core/benches/verdict_construction.rs.
% Raw sweep data: data/sweep_raw_20260914T171258Z.csv (env A),
% data/sweep_raw_20260914T164735Z_runnervmlun5p.csv (env B), and
% data/sweep_raw_20260914T231541Z.csv (env C; fingerprints in data/sweep_env_*.txt).

\section{Empirical Validation}
\label{sec:benchmarks}

The Constitutional Enclosure Theorem guarantees that silent decisions
are compile-time errors.  A natural question follows: does the
compile-time guarantee impose a \emph{runtime cost}?  This section
answers that question across two independent hardware platforms, under
concurrency, and against two different readings of "the status quo"
rather than one hand-picked baseline.

% ----------------------------------------------------------------
\subsection{Experimental Setup}
\label{sec:setup}

\paragraph{Hardware and software.}
Table~\ref{tab:crossplatform} compares two platforms of different
vendor and model: (A) Intel Xeon (Model 173), 2~vCPU, KVM; (B) AMD EPYC
9V74, 4~vCPU (Azure). Table~\ref{tab:fivemode}'s five-mode contrast
uses a third, more heavily shared/virtualized container (C), whose
absolute latencies are not comparable to A/B (caveat below the table).
All three run \texttt{rustc}~1.94.1, pinned by
\texttt{rust-toolchain.toml}. Criterion~0.5.1 produces
Table~\ref{tab:criterion}; a dedicated harness
(\texttt{btv-core/benches/sweep\_concurrent.rs}) produces
Tables~\ref{tab:crossplatform}--\ref{tab:fivemode}.

\paragraph{Sweep methodology.}
Each (payload, mode, thread-count) configuration runs 5 trials; each
Rayon worker times its own operations with a thread-local
P\textsuperscript{2} quantile estimator~\cite{jain1985p2} ($O(1)$
memory, no stored observations). \textbf{Every p95/p99 figure in this
section is therefore a P\textsuperscript{2} estimate, not an exact
order statistic.} No harness-owned lock sits in the timed path, except
\texttt{FullPipelineDurable} (Table~\ref{tab:fivemode}), which
deliberately shares one \texttt{SqliteLogSink} because
\texttt{SQLite}'s own locking is what is being measured there. The
committed runs use a reduced wall-clock target (1--2\,s/config, not the
project's 90\,s standard); a dedicated-hardware 90\,s run is still owed
but would not change the ratios below, which compare modes measured
together on the same run, not absolute numbers against other papers.

\paragraph{Memory footprint.}
\begin{center}
\texttt{std::mem::size\_of::\textless{}Verdict\textgreater{}() = 184 bytes} (x86-64, rustc~1.94.1)
\end{center}
32-byte \texttt{Blake3Hash}; \texttt{ComplianceToken} (two strings, a
\texttt{u32}, and a 32-byte authority signature); \texttt{Decision}
discriminant; explanation \texttt{String}; 32-byte integrity seal.

% ----------------------------------------------------------------
\subsection{Results}
\label{sec:results}

\paragraph{Cross-platform latency and variance.}
Table~\ref{tab:crossplatform}: \texttt{issue\_verdict} end-to-end
(BLAKE3 + HMAC + in-memory append), 1 thread, median and CV over 5
trials.

\begin{table}[t]
\caption{Full pipeline latency by payload and platform (1 thread, median of 5 trials)}
\label{tab:crossplatform}
\centering\small
\begin{tabular}{lrrrr}
\hline
\textbf{Payload} & \textbf{A: mean (ns)} & \textbf{A: CV} & \textbf{B: mean (ns)} & \textbf{B: CV} \\
\hline
64~B   & 383   & 0.7\% & 411   & 0.5\% \\
512~B  & 746   & 0.3\% & 889   & 0.2\% \\
4~KiB  & 1294  & 0.3\% & 1735  & 0.2\% \\
\hline
\end{tabular}
\end{table}

CV under 1\% at every payload on both platforms --- repeatability
across separate runs, a stronger statement than one run's bootstrap CI.

\paragraph{Thread scaling.}
On B, \texttt{full\_pipeline} throughput at 64~B: 1.96~M~ops/s (1
thread) $\to$ 3.73~M (2) $\to$ 5.70~M (4) --- within 3\% of ideal
linear scaling, consistent with no structural lock in the measured
path.

\paragraph{Five-mode contrast.}
Table~\ref{tab:fivemode} adds what a single baseline omitted: a
durable arm (\texttt{FullPipelineDurable}: on-disk \texttt{SQLite},
WAL + \texttt{synchronous=FULL} actually in effect, not an in-memory
sink that silently ignores both) and a status-quo arm symmetric with
BTV (\texttt{StatusQuoDigestLog} logs a digest, like
\texttt{FullPipeline}, instead of the full context hex-encoded). 4~KiB
payload, 4 threads, environment C:

\begin{table}[t]
\caption{Five-mode contrast, 4~KiB payload, 4 threads, environment C (median of 5 trials)}
\label{tab:fivemode}
\centering\small
\begin{tabular}{lrr}
\hline
\textbf{Mode} & \textbf{Mean latency} & \textbf{Throughput (ops/s)} \\
\hline
\texttt{verdict\_only}            & 4.34~\textmu{}s   & 505{,}937 \\
\texttt{status\_quo\_digest\_log} & 1.53~\textmu{}s   & 2{,}268{,}330 \\
\texttt{full\_pipeline}           & 7.27~\textmu{}s   & 511{,}108 \\
\texttt{status\_quo\_async\_log}  & 19.0~\textmu{}s   & 196{,}186 \\
\texttt{full\_pipeline\_durable}  & 168.6~\textmu{}s  & 22{,}338 \\
\hline
\end{tabular}
\end{table}

A sensitivity range, not one number: against a full-context logger,
BTV is 2.6$\times$ faster at 4~KiB; against a digest-only logger,
4.7$\times$ slower --- a modeling choice about what "the status quo"
logs, not a property of BTV, so both are shown. Durability is the
dominant cost, by more than an order of magnitude: 23$\times$ slower
than the in-memory pipeline at the same payload and thread count. One
anomaly is reported rather than smoothed over: at 1 thread,
\texttt{full\_pipeline\_durable}'s CV exceeds 140\% (vs.\ 2--3\% at 2
and 4 threads), plausibly a single-threaded \texttt{SQLite} WAL
file's page-cache warm-up effect --- flagged as unexplained, not
presented as stable. \emph{Environment C's absolute latencies (higher
than Table~\ref{tab:crossplatform}) are not comparable across tables;
only ratios within this table, from the same run, are load-bearing.}

\paragraph{Single-operation cost breakdown.}
Table~\ref{tab:criterion} isolates construction from persistence
(Criterion, 95\% CI, 100 samples, environment A, fixed small context).

\begin{table}[t]
\caption{Single-operation latency (Criterion, 95\% CI)}
\label{tab:criterion}
\centering\small
\begin{tabular}{lr}
\hline
\textbf{Operation} & \textbf{Mean} \\
\hline
\texttt{Verdict::new} (construction only)          & 5.21~\textmu{}s \\
\texttt{issue\_verdict}, in-memory sink            & 6.64~\textmu{}s \\
\texttt{issue\_verdict}, durable \texttt{SQLite}   & 1.28~ms \\
\hline
\end{tabular}
\end{table}

The jump to durable persistence ($\approx$246$\times$) is the real
price of non-repudiation once "durable" means an \texttt{fsync}, not a
RAM insert; the type-level discipline contributes none of it, since
private fields, \texttt{\#[must\_use]}, and \texttt{pub(crate)}
constructors emit no code.

% ----------------------------------------------------------------
\subsection{Discussion}
\label{sec:bench-discussion}

\paragraph{What the type system costs, and what it does not.}
The linear-type discipline is compile-time only: it adds no branch,
allocation, or synchronisation primitive, and Table~\ref{tab:criterion}
is consistent with that --- the added cost is exactly the cryptography
and I/O invoked, nothing more. That is narrower than "the guarantee is
free": Table~\ref{tab:fivemode} shows it costs one to four orders of
magnitude more than doing nothing, depending on what "nothing" means.
What the theorem buys is that the expensive path (durable, signed,
non-repudiable) is the \emph{only} path the type system allows for a
materialised \texttt{Verdict} --- there is no cheaper, silently
unaccountable one available by omission.

\paragraph{Verification cost.}
\texttt{Verdict::verify\_integrity} recomputes an HMAC over fixed-size
fields without re-hashing the original context ($\ll$1~\textmu{}s,
subsumed in \texttt{verdict\_only} above): retroactive audit at
population scale is computationally unremarkable, independent of the
persistence cost above.

% ----------------------------------------------------------------
\subsection{Worked Examples}
\label{sec:worked-examples}

Four worked examples demonstrate BTV in semi-realistic scenarios: credit
scoring under LGPD (30-day appeal window), hiring screening under EU AI Act
Art.~86 (high-risk classification), emergency triage (where \texttt{Allow}
decisions also require evidence), and emergency escalation to a human
operator when automated evidence cannot be produced (Corollary~\ref{cor:escalated}).
Each constructs verdicts from domain-specific payloads and verifies
integrity; source is in the companion repository (\texttt{examples/}).
